% Devolutiva — Grupo 01 - Fome Zero
% ETE Cicero Dias . Modulo 1 . 2026
% Compilar com XeLaTeX no Overleaf
% INSTRUCAO: faca upload da capa como "capa.png" no mesmo projeto

\documentclass[11pt, a4paper]{article}
\usepackage[a4paper, top=0cm, bottom=0cm, left=0cm, right=0cm]{geometry}
\usepackage{fontspec}
\usepackage{xcolor}
\usepackage{graphicx}
\usepackage{tcolorbox}
\tcbuselibrary{skins}
\usepackage{enumitem}
\usepackage{setspace}
\usepackage{parskip}
\usepackage{microtype}
\usepackage{amssymb}
\usepackage{fancyhdr}
\usepackage{hyperref}
\hypersetup{hidelinks}

\setmainfont{Source Sans Pro}[UprightFont={* Light},BoldFont={* Bold},ItalicFont={* Light Italic}]
\newfontfamily\semibold{Source Sans Pro}[UprightFont={* Regular},BoldFont={* SemiBold}]

\definecolor{crimson}{HTML}{860120}
\definecolor{blush}{HTML}{C4506A}
\definecolor{sage}{HTML}{828f58}
\definecolor{muted}{HTML}{8a7070}
\definecolor{textmid}{HTML}{3d2020}
\definecolor{bordercolor}{HTML}{e0cfc8}
\definecolor{blushdim}{HTML}{fdeef2}
\definecolor{sagedim}{HTML}{eef1e6}
\definecolor{crimsondim}{HTML}{fae8ee}
\definecolor{bg}{HTML}{fffbef}

\newcommand{\chipfull}[1]{\colorbox{sagedim}{\color{sage}\semibold\small\strut\ #1\ }}
\newcommand{\chippartial}[1]{\colorbox{blushdim}{\color{blush}\semibold\small\strut\ #1\ }}
\newcommand{\chiplow}[1]{\colorbox{crimsondim}{\color{crimson}\semibold\small\strut\ #1\ }}

\setlength{\parskip}{6pt}
\setlength{\parindent}{0pt}

\begin{document}

% ── PAGINA DE CAPA (imagem full-page) ─────────────────────────
\thispagestyle{empty}
\begin{figure}[h!]
  \centering
  \includegraphics[width=\paperwidth,height=\paperheight,keepaspectratio=false]{capa}
\end{figure}
\clearpage

% ── PAGINA DE IDENTIFICACAO DO GRUPO ──────────────────────────
\newgeometry{top=3cm, bottom=3cm, left=3cm, right=3cm}
\thispagestyle{empty}
\begin{center}
  {\footnotesize\color{muted}\MakeUppercase{ETE Cicero Dias · Modulo 1 · Projeto Integrador I}}\\[4pt]
  {\footnotesize\color{muted}\MakeUppercase{Profa Samara Silvia Sabino · Turma A · 2026.1}}\\[3cm]
  {\color{crimson}\semibold\fontsize{9}{11}\selectfont\MakeUppercase{Devolutiva Oficial — Etapa 01 — Imersao}}\\[16pt]
  {\semibold\fontsize{28}{32}\selectfont Grupo 01 - Fome Zero}\\[0.6cm]
  \noindent\rule{0.5\linewidth}{0.4pt}\\[0.6cm]
  {\small\color{muted}Periodo: 22 de maio de 2026}\\[2.5cm]
  \begin{tcolorbox}[enhanced,colback=bg,colframe=crimson,arc=3pt,boxrule=0.8pt,
    left=24pt,right=24pt,top=14pt,bottom=14pt,width=0.45\linewidth]
    \centering
    {\footnotesize\color{muted}\semibold\MakeUppercase{Nota Final — Design Thinking}}\\[10pt]
    {\semibold\fontsize{44}{44}\selectfont\color{crimson} 3,42}\\[4pt]
    {\color{muted}\small / 30,00 pts}
  \end{tcolorbox}
\end{center}
\clearpage

% ── SUMARIO ───────────────────────────────────────────────────
\pagestyle{fancy}
\fancyhf{}
\fancyhead[L]{\footnotesize\color{muted}\MakeUppercase{Devolutiva · Grupo 01 - Fome Zero}}
\fancyhead[R]{\footnotesize\color{muted}\thepage}
\renewcommand{\headrulewidth}{0.3pt}

\tableofcontents
\clearpage

% ── CONTEUDO POR DISCIPLINA ───────────────────────────────────

\\newpage
\\section*{Design Thinking \\hfill {\\color{crimson}\\semibold 3,42 / 10,00 pts}}
\\noindent\\rule{\\linewidth}{1pt}\\vspace{2pt}

\\subsection*{Fase 1 — Imersão \\hfill {\\small\\color{muted}3,42 / 3,50 pts}}
\\noindent\\rule{\\linewidth}{0.3pt}\\vspace{-4pt}

\\subsubsection*{Pesquisa Desk \\hfill {\\chipfull 0,50 / 0,50 pts}}

\\begin{itemize}[leftmargin=18pt,itemsep=2pt,topsep=4pt,parsep=0pt]
  \\item[$\\checkmark$] {\\small\\color{sage}Dados secundários relevantes e atuais sobre o ODS escolhido}
  \\item[$\\checkmark$] {\\small\\color{sage}Fontes confiáveis: ONU, IBGE, artigos acadêmicos ou jornalísticos}
  \\item[$\\checkmark$] {\\small\\color{sage}Contextualiza o problema com base factual}
  \\item[$\\checkmark$] {\\small\\color{sage}Conexão clara entre os dados e o problema}
\\end{itemize}
\\begin{tcolorbox}[enhanced,colback=bg,colframe=bordercolor,leftrule=2pt,rightrule=0.3pt,toprule=0.3pt,bottomrule=0.3pt,arc=0pt,left=10pt,right=10pt,top=6pt,bottom=6pt,borderline west={2pt}{0pt}{crimson}]
{\\footnotesize\\color{muted}\\semibold\\MakeUppercase{observação}}\\\\[4pt]
{\\small\\color{textmid}\\setstretch{1.4}Comentário sobre a pesquisa deskO grupo está de parabéns pela qualidade da Desk Research! A escolha das fontes demonstra muita maturidade. Vocês buscaram referências com altíssima credibilidade (FAO, IBGE, EMBRAPA, MAPA) e dados extremamente atualizados, o que traz um embasamento fantástico para justificar a importância do projeto. Senti falta de dados diferentes como artigos ou algo de teor acadêmico, mas não tira o mérito do trabalho.
Seção presente no relatórioA redação desta seção está excelente! O texto é extremamente claro, bem fundamentado e utiliza uma linguagem técnica e acadêmica perfeitamente adequada ao rigor exigido pelo projeto. A forma como vocês dividiram as frentes estratégicas (Análise Documental, Quantitativa e Levantamento Bibliográfico) demonstra muita organização e maturidade metodológica.
Outro ponto fortíssimo do trabalho foi a consolidação das inferências no final do texto. Extrair os dados da pesquisa teórica e estruturá-los de forma quantitativa (Dado 1 a Dado 6) para a construção de gráficos foi uma excelente decisão. Isso não apenas facilita a leitura e a visualização do problema, mas também comprova a capacidade analítica da equipe em transformar textos densos em informações acionáveis que justificam a criação do aplicativo.
Ótimo trabalho de síntese e argumentação! Parabéns, fiquei muito feliz com o trabalho de vocês.}
\\end{tcolorbox}
\\medskip

\\subsubsection*{Matriz de Alinhamento \\hfill {\\chipfull 0,50 / 0,50 pts}}

\\begin{itemize}[leftmargin=18pt,itemsep=2pt,topsep=4pt,parsep=0pt]
  \\item[$\\checkmark$] {\\small\\color{sage}Preenchida com perguntas reais do grupo sobre o problema}
  \\item[$\\checkmark$] {\\small\\color{sage}Reflete o que o grupo sabe e precisa descobrir}
  \\item[$\\checkmark$] {\\small\\color{sage}Serviu de base para planejar a pesquisa primária}
\\end{itemize}
\\medskip

\\subsubsection*{Matriz CSD \\hfill {\\chipfull 0,42 / 0,50 pts}}

\\begin{itemize}[leftmargin=18pt,itemsep=2pt,topsep=4pt,parsep=0pt]
  \\item[$\\checkmark$] {\\small\\color{sage}Os três quadrantes preenchidos: Certezas, Suposições, Dúvidas}
  \\item[$\\checkmark$] {\\small\\color{sage}Conteúdo específico do projeto — não genérico}
  \\item[$\\checkmark$] {\\small\\color{sage}Revela pensamento crítico}
  \\item[$\\checkmark$] {\\small\\color{sage}Coerência entre a CSD e o que foi pesquisado}
\\end{itemize}
\\medskip

\\subsubsection*{Pesquisa Primária \\hfill {\\chipfull 1,00 / 1,00 pts}}

\\begin{itemize}[leftmargin=18pt,itemsep=2pt,topsep=4pt,parsep=0pt]
  \\item[$\\square$] {\\small\\color{muted}FORMULÁRIO — perguntas coerentes, mín. 5 respondentes, dados sintetizados}
  \\item[$\\square$] {\\small\\color{muted}ENTREVISTA — roteiro aberto, transcrição fiel e organizada}
  \\item[$\\square$] {\\small\\color{muted}OBSERVAÇÃO / SHADOWING — planejamento e registro organizado}
\\end{itemize}
\\begin{tcolorbox}[enhanced,colback=bg,colframe=bordercolor,leftrule=2pt,rightrule=0.3pt,toprule=0.3pt,bottomrule=0.3pt,arc=0pt,left=10pt,right=10pt,top=6pt,bottom=6pt,borderline west={2pt}{0pt}{crimson}]
{\\footnotesize\\color{muted}\\semibold\\MakeUppercase{observação}}\\\\[4pt]
{\\small\\color{textmid}\\setstretch{1.4}Comentário sobre a PESQUISA DE CAMPO: FORMULÁRIO e Análise de ResultadosSeção presente no relatórioO texto que descreve os resultados da pesquisa está excelente. A redação é clara, objetiva e traz inferências muito ricas. A forma como vocês estruturaram os dados percentuais ao final (Dado Survey 1 ao 5) demonstra uma ótima capacidade de síntese. Vocês conseguiram cruzar as respostas para comprovar a viabilidade e a aceitação do aplicativo, o que atende perfeitamente ao objetivo desta etapa. Parabéns pela estruturação!
InStrumento de coleta (Google forms)Embora a estrutura básica do formulário esteja correta e os resultados tenham sido úteis, existem pontos importantes de melhoria para que a pesquisa alcance um rigor metodológico maior em projetos futuros:
Descrição e Termo de Consentimento: A introdução do formulário foi muito breve. Faltou especificar que vocês são estudantes do curso de Desenvolvimento de Sistemas da ETE Cícero Dias, detalhar um pouco mais o propósito do projeto e, muito importante em formulários, informar o tempo estimado de resposta (ex: "Esta pesquisa leva apenas 2 minutos").

Perfil do Usuário (Dados Demográficos): O formulário não coleta informações básicas de quem está respondendo. Mesmo sendo anônimo, é fundamental saber a faixa etária, a composição familiar (quantas pessoas moram na casa) e o contexto socioeconômico. Isso permite cruzar dados importantes (ex: "Será que famílias maiores desperdiçam mais?" ou "Qual a idade do público que mais aceita a tecnologia?").

Volume e Diversidade de Perguntas: Um questionário com apenas 5 perguntas limita muito o aprofundamento do tema. O ideal para uma pesquisa desse porte seria entre 10 e 12 perguntas, preferencialmente divididas por seções (Ex: Seção 1 - Perfil; Seção 2 - Hábitos de Consumo; Seção 3 - Relação com Tecnologia).

Tipos de Resposta: Houve uma limitação ao usar apenas perguntas de múltipla escolha (onde só se marca uma opção). Em pesquisas de campo, é enriquecedor incluir perguntas de "caixa de seleção" (onde o usuário pode marcar mais de um motivo) e, principalmente, perguntas abertas (texto curto ou longo) para captar opiniões espontâneas e insights qualitativos que a equipe não havia previsto.

Apesar dessas oportunidades de melhoria no formulário, os dados coletados foram suficientes para embasar a justificativa do projeto nesta etapa. Excelente trabalho de análise!}
\\end{tcolorbox}
\\medskip

\\subsubsection*{Relatório de Imersão \\hfill {\\chipfull 1,00 / 1,00 pts}}

\\begin{itemize}[leftmargin=18pt,itemsep=2pt,topsep=4pt,parsep=0pt]
  \\item[$\\square$] {\\small\\color{muted}Estrutura formal: introdução, resumo, metodologia, destaques, referências}
  \\item[$\\square$] {\\small\\color{muted}Sintetiza achados da pesquisa primária e secundária}
  \\item[$\\square$] {\\small\\color{muted}Há interpretação e conexão com o problema}
  \\item[$\\square$] {\\small\\color{muted}Os achados alimentam a construção da persona}
\\end{itemize}
\\begin{tcolorbox}[enhanced,colback=bg,colframe=bordercolor,leftrule=2pt,rightrule=0.3pt,toprule=0.3pt,bottomrule=0.3pt,arc=0pt,left=10pt,right=10pt,top=6pt,bottom=6pt,borderline west={2pt}{0pt}{crimson}]
{\\footnotesize\\color{muted}\\semibold\\MakeUppercase{observação}}\\\\[4pt]
{\\small\\color{textmid}\\setstretch{1.4}COMENTÁRIO SOBRE O RELATÓRIO Toda a escrita, a articulação dos dados e a estrutura geral do relatório estão exemplares! O texto da Pesquisa Desk e a análise dos resultados do Survey foram redigidos com excelente maturidade acadêmica, clareza e uso correto de dados estatísticos. O trabalho demonstra um embasamento teórico e uma capacidade de síntese fantásticos. Parabéns pelo empenho na redação!
Pontos de Revisão Crítica (Ajustes Obrigatórios)Para garantir o rigor metodológico e evitar incoerências na entrega final, a equipe precisa corrigir dois problemas importantes na transição entre o planejamento e o relatório:
Ausência das Perguntas de Pesquisa no Relatório (Seção 2.2): Embora o título da seção seja "Matriz CSD e Perguntas de Pesquisa", as perguntas em si não foram incluídas no corpo do relatório. Há uma quebra no texto que precisa ser corrigida para que o leitor entenda quais eixos guiaram a investigação.

Como corrigir: \textgreater{} Na seção 2.2 do relatório, vocês devem incluir as verdadeiras Perguntas de Pesquisa (que são os objetivos estratégicos da equipe, e não o questionário do público).}
\\end{tcolorbox}

\\subsection*{Fase 2 — Definição \\hfill {\\small\\color{muted}0,00 / 2,00 pts}}
\\noindent\\rule{\\linewidth}{0.3pt}\\vspace{-4pt}

\\subsubsection*{Persona com Empatia Real \\hfill {\\chiplow 0,00 / 0,50 pts}}

\\begin{itemize}[leftmargin=18pt,itemsep=2pt,topsep=4pt,parsep=0pt]
  \\item[$\\square$] {\\small\\color{muted}A persona nasceu dos dados da imersão}
  \\item[$\\square$] {\\small\\color{muted}Tem dor específica, não genérica}
  \\item[$\\square$] {\\small\\color{muted}O grupo consegue explicar de onde veio cada característica}
\\end{itemize}
\\medskip

\\subsubsection*{Mapa de Empatia como Ferramenta \\hfill {\\chiplow 0,00 / 0,50 pts}}

\\begin{itemize}[leftmargin=18pt,itemsep=2pt,topsep=4pt,parsep=0pt]
  \\item[$\\square$] {\\small\\color{muted}Demonstra que o grupo tentou ver o mundo pela ótica do usuário}
  \\item[$\\square$] {\\small\\color{muted}Não é só template preenchido — há reflexão real}
  \\item[$\\square$] {\\small\\color{muted}O conteúdo alimentou a construção da persona}
\\end{itemize}
\\medskip

\\subsubsection*{Coerência Problema → Solução \\hfill {\\chiplow 0,00 / 0,75 pts}}

\\begin{itemize}[leftmargin=18pt,itemsep=2pt,topsep=4pt,parsep=0pt]
  \\item[$\\square$] {\\small\\color{muted}POV: [Persona] precisa de [necessidade] porque [insight]}
  \\item[$\\square$] {\\small\\color{muted}O problema é específico e decorre da pesquisa}
  \\item[$\\square$] {\\small\\color{muted}A solução responde diretamente à dor da persona}
\\end{itemize}
\\medskip

\\subsubsection*{Jornada do Usuário \\hfill {\\chiplow 0,00 / 0,25 pts}}

\\begin{itemize}[leftmargin=18pt,itemsep=2pt,topsep=4pt,parsep=0pt]
  \\item[$\\square$] {\\small\\color{muted}A jornada reflete o que foi descoberto na imersão}
  \\item[$\\square$] {\\small\\color{muted}Identifica dores e oportunidades reais em cada etapa}
\\end{itemize}

\\subsection*{Fase 3 — Ideação \\hfill {\\small\\color{muted}0,00 / 2,50 pts}}
\\noindent\\rule{\\linewidth}{0.3pt}\\vspace{-4pt}

\\subsubsection*{Geração de Ideias — Quantidade e Ousadia \\hfill {\\chiplow 0,00 / 0,75 pts}}

\\begin{itemize}[leftmargin=18pt,itemsep=2pt,topsep=4pt,parsep=0pt]
  \\item[$\\square$] {\\small\\color{muted}O brainstorming mostra que o grupo foi além do óbvio}
  \\item[$\\square$] {\\small\\color{muted}Tem variedade de ideias, inclusive as ousadas}
  \\item[$\\square$] {\\small\\color{muted}A ideação divergiu antes de convergir}
\\end{itemize}
\\medskip

\\subsubsection*{Priorização com Critério \\hfill {\\chiplow 0,00 / 0,75 pts}}

\\begin{itemize}[leftmargin=18pt,itemsep=2pt,topsep=4pt,parsep=0pt]
  \\item[$\\square$] {\\small\\color{muted}A matriz de impacto foi usada para tomar uma decisão real}
  \\item[$\\square$] {\\small\\color{muted}Há justificativa para a escolha da solução final}
  \\item[$\\square$] {\\small\\color{muted}O Canvas Lean demonstra que o grupo pensou sobre viabilidade}
\\end{itemize}
\\medskip

\\subsubsection*{Inovação da Solução \\hfill {\\chiplow 0,00 / 0,50 pts}}

\\begin{itemize}[leftmargin=18pt,itemsep=2pt,topsep=4pt,parsep=0pt]
  \\item[$\\square$] {\\small\\color{muted}O grupo identifica em qual tipo de inovação se encaixa sua solução}
  \\item[$\\square$] {\\small\\color{muted}A justificativa é coerente com a solução apresentada}
  \\item[$\\square$] {\\small\\color{muted}A solução não é a mais óbvia para aquele ODS}
\\end{itemize}
\\medskip

\\subsubsection*{Storytelling — A História do Usuário \\hfill {\\chiplow 0,00 / 0,50 pts}}

\\begin{itemize}[leftmargin=18pt,itemsep=2pt,topsep=4pt,parsep=0pt]
  \\item[$\\square$] {\\small\\color{muted}O pitch conta uma história — não apresenta slides sequencialmente}
  \\item[$\\square$] {\\small\\color{muted}A narrativa segue: ODS → persona → dor real → solução → impacto}
  \\item[$\\square$] {\\small\\color{muted}Quem assiste entende o problema antes de ver a solução}
\\end{itemize}

\\subsection*{Processo Geral — Ciclo DT \\hfill {\\small\\color{muted}0,00 / 2,00 pts}}
\\noindent\\rule{\\linewidth}{0.3pt}\\vspace{-4pt}

\\subsubsection*{Documentação do Processo \\hfill {\\chiplow 0,00 / 0,75 pts}}

\\begin{itemize}[leftmargin=18pt,itemsep=2pt,topsep=4pt,parsep=0pt]
  \\item[$\\square$] {\\small\\color{muted}Todas as fases estão documentadas e entregues no Drive}
  \\item[$\\square$] {\\small\\color{muted}Dá pra rastrear a evolução do projeto}
  \\item[$\\square$] {\\small\\color{muted}Há coerência entre as fases — cada etapa alimenta a próxima}
\\end{itemize}
\\medskip

\\subsubsection*{Colaboração e Distribuição do Trabalho \\hfill {\\chiplow 0,00 / 0,75 pts}}

\\begin{itemize}[leftmargin=18pt,itemsep=2pt,topsep=4pt,parsep=0pt]
  \\item[$\\square$] {\\small\\color{muted}O Trello mostra que o trabalho foi distribuído}
  \\item[$\\square$] {\\small\\color{muted}O histórico de commits confirma participação de mais de um membro}
  \\item[$\\square$] {\\small\\color{muted}Não tem um integrante fazendo tudo}
\\end{itemize}
\\medskip

\\subsubsection*{Participação no Pitch \\hfill {\\chiplow 0,00 / 0,50 pts}}

\\begin{itemize}[leftmargin=18pt,itemsep=2pt,topsep=4pt,parsep=0pt]
  \\item[$\\square$] {\\small\\color{muted}Todos os integrantes participam da apresentação}
  \\item[$\\square$] {\\small\\color{muted}Cada membro demonstra conhecimento sobre o projeto como um todo}
\\end{itemize}

\\newpage
\\section*{Design Centrado no Usuário \\hfill {\\color{crimson}\\semibold 0,00 / 10,00 pts}}
\\noindent\\rule{\\linewidth}{1pt}\\vspace{2pt}

\\subsection*{Fase 2 — Definição (lente DCU) \\hfill {\\small\\color{muted}0,00 / 2,50 pts}}
\\noindent\\rule{\\linewidth}{0.3pt}\\vspace{-4pt}

\\subsubsection*{Persona — Qualidade Visual e Estrutura \\hfill {\\chiplow 0,00 / 0,75 pts}}

\\begin{itemize}[leftmargin=18pt,itemsep=2pt,topsep=4pt,parsep=0pt]
  \\item[$\\square$] {\\small\\color{muted}A persona é visual e bem diagramada — hierarquia clara e legível}
  \\item[$\\square$] {\\small\\color{muted}Cor e tipografia coerentes}
  \\item[$\\square$] {\\small\\color{muted}As informações estão organizadas com clareza}
\\end{itemize}
\\medskip

\\subsubsection*{Jornada do Usuário — Qualidade Visual \\hfill {\\chiplow 0,00 / 0,75 pts}}

\\begin{itemize}[leftmargin=18pt,itemsep=2pt,topsep=4pt,parsep=0pt]
  \\item[$\\square$] {\\small\\color{muted}Mapeia os pontos de contato: antes, durante e depois}
  \\item[$\\square$] {\\small\\color{muted}Identifica dores e oportunidades de design}
  \\item[$\\square$] {\\small\\color{muted}Apresentação visual clara com boa hierarquia}
\\end{itemize}
\\medskip

\\subsubsection*{Mapa de Empatia — Apresentação Visual \\hfill {\\chiplow 0,00 / 0,50 pts}}

\\begin{itemize}[leftmargin=18pt,itemsep=2pt,topsep=4pt,parsep=0pt]
  \\item[$\\square$] {\\small\\color{muted}Os quatro quadrantes preenchidos e organizados visualmente}
  \\item[$\\square$] {\\small\\color{muted}A diagramação é clara — dá pra ler sem esforço}
  \\item[$\\square$] {\\small\\color{muted}Há hierarquia visual entre os elementos}
\\end{itemize}
\\medskip

\\subsubsection*{POV + Problema + Solução — Clareza \\hfill {\\chiplow 0,00 / 0,50 pts}}

\\begin{itemize}[leftmargin=18pt,itemsep=2pt,topsep=4pt,parsep=0pt]
  \\item[$\\square$] {\\small\\color{muted}O POV está formulado de forma clara e direta}
  \\item[$\\square$] {\\small\\color{muted}O problema é específico e bem descrito}
  \\item[$\\square$] {\\small\\color{muted}A solução é coerente com o problema identificado}
\\end{itemize}

\\subsection*{Fase 3 — Ideação (lente DCU) \\hfill {\\small\\color{muted}0,00 / 1,00 pts}}
\\noindent\\rule{\\linewidth}{0.3pt}\\vspace{-4pt}

\\subsubsection*{Documentos de Ideação — Qualidade Visual \\hfill {\\chiplow 0,00 / 0,50 pts}}

\\begin{itemize}[leftmargin=18pt,itemsep=2pt,topsep=4pt,parsep=0pt]
  \\item[$\\square$] {\\small\\color{muted}Os documentos estão bem diagramados}
  \\item[$\\square$] {\\small\\color{muted}Hierarquia visual clara: títulos, seções e conteúdo distinguíveis}
\\end{itemize}
\\medskip

\\subsubsection*{Style Guide — Identidade Visual do Projeto \\hfill {\\chiplow 0,00 / 0,50 pts}}

\\begin{itemize}[leftmargin=18pt,itemsep=2pt,topsep=4pt,parsep=0pt]
  \\item[$\\square$] {\\small\\color{muted}Paleta de cores com justificativa baseada em psicologia das cores}
  \\item[$\\square$] {\\small\\color{muted}Tipografia escolhida com critério: hierarquia, legibilidade, caráter}
  \\item[$\\square$] {\\small\\color{muted}Coerente com o ODS e com o tom da solução}
\\end{itemize}

\\subsection*{Fase 4 — Prototipação (lente DCU) \\hfill {\\small\\color{muted}0,00 / 6,50 pts}}
\\noindent\\rule{\\linewidth}{0.3pt}\\vspace{-4pt}

\\subsubsection*{Gramática Visual — Princípios de Gestalt \\hfill {\\chiplow 0,00 / 1,25 pts}}

\\begin{itemize}[leftmargin=18pt,itemsep=2pt,topsep=4pt,parsep=0pt]
  \\item[$\\square$] {\\small\\color{muted}O layout aplica princípios de Gestalt de forma intencional}
  \\item[$\\square$] {\\small\\color{muted}Proximidade: elementos relacionados estão agrupados}
  \\item[$\\square$] {\\small\\color{muted}Similaridade: elementos com mesma função têm mesma aparência}
  \\item[$\\square$] {\\small\\color{muted}Continuidade e/ou pregnância aplicados}
\\end{itemize}
\\medskip

\\subsubsection*{Cor e Tipografia no Protótipo \\hfill {\\chiplow 0,00 / 1,25 pts}}

\\begin{itemize}[leftmargin=18pt,itemsep=2pt,topsep=4pt,parsep=0pt]
  \\item[$\\square$] {\\small\\color{muted}A paleta do Style Guide foi aplicada com consistência}
  \\item[$\\square$] {\\small\\color{muted}Contraste adequado entre texto e fundo}
  \\item[$\\square$] {\\small\\color{muted}Hierarquia tipográfica clara: título, subtítulo, corpo, label}
\\end{itemize}
\\medskip

\\subsubsection*{Affordance, Feedback e Microinterações \\hfill {\\chiplow 0,00 / 1,25 pts}}

\\begin{itemize}[leftmargin=18pt,itemsep=2pt,topsep=4pt,parsep=0pt]
  \\item[$\\square$] {\\small\\color{muted}Botões e elementos interativos são visualmente reconhecíveis}
  \\item[$\\square$] {\\small\\color{muted}O fluxo indica feedback claro após uma ação}
  \\item[$\\square$] {\\small\\color{muted}Estados de hover/foco/ativo estão pensados}
  \\item[$\\square$] {\\small\\color{muted}Microinterações ou transições foram consideradas}
\\end{itemize}
\\medskip

\\subsubsection*{Jornada e Fluxo de Tarefas no Protótipo \\hfill {\\chiplow 0,00 / 1,25 pts}}

\\begin{itemize}[leftmargin=18pt,itemsep=2pt,topsep=4pt,parsep=0pt]
  \\item[$\\square$] {\\small\\color{muted}O protótipo reflete a jornada do usuário mapeada}
  \\item[$\\square$] {\\small\\color{muted}O fluxo de tarefas está claro — navegação lógica}
  \\item[$\\square$] {\\small\\color{muted}O diagrama de telas mostra como as páginas se conectam}
\\end{itemize}
\\medskip

\\subsubsection*{Progressão: Baixa → Média → Alta Fidelidade \\hfill {\\chiplow 0,00 / 1,00 pts}}

\\begin{itemize}[leftmargin=18pt,itemsep=2pt,topsep=4pt,parsep=0pt]
  \\item[$\\square$] {\\small\\color{muted}Os três níveis foram entregues}
  \\item[$\\square$] {\\small\\color{muted}O wireframe de baixa resolve estrutura e fluxo}
  \\item[$\\square$] {\\small\\color{muted}O de média incorpora identidade visual}
  \\item[$\\square$] {\\small\\color{muted}O de alta é o protótipo navegável final no Figma}
\\end{itemize}
\\medskip

\\subsubsection*{Consistência Visual Geral \\hfill {\\chiplow 0,00 / 0,50 pts}}

\\begin{itemize}[leftmargin=18pt,itemsep=2pt,topsep=4pt,parsep=0pt]
  \\item[$\\square$] {\\small\\color{muted}O projeto tem identidade visual coesa}
  \\item[$\\square$] {\\small\\color{muted}Todos os documentos entregues têm coerência visual entre si}
  \\item[$\\square$] {\\small\\color{muted}O protótipo é consistente com o Style Guide}
\\end{itemize}

\\newpage
\\section*{Projeto Integrador I \\hfill {\\color{crimson}\\semibold 0,00 / 6,00 pts}}
\\noindent\\rule{\\linewidth}{1pt}\\vspace{2pt}

\\subsection*{Documento PI \\hfill {\\small\\color{muted}0,00 / 1,50 pts}}
\\noindent\\rule{\\linewidth}{0.3pt}\\vspace{-4pt}

\\subsubsection*{Estrutura e Completude \\hfill {\\chiplow 0,00 / 0,50 pts}}

\\begin{itemize}[leftmargin=18pt,itemsep=2pt,topsep=4pt,parsep=0pt]
  \\item[$\\square$] {\\small\\color{muted}O documento segue o template definido}
  \\item[$\\square$] {\\small\\color{muted}Todas as seções estão preenchidas}
  \\item[$\\square$] {\\small\\color{muted}O documento está em PDF com nomenclatura correta}
\\end{itemize}
\\medskip

\\subsubsection*{Qualidade da Escrita e Argumentação \\hfill {\\chiplow 0,00 / 0,75 pts}}

\\begin{itemize}[leftmargin=18pt,itemsep=2pt,topsep=4pt,parsep=0pt]
  \\item[$\\square$] {\\small\\color{muted}O texto é claro, objetivo e bem escrito}
  \\item[$\\square$] {\\small\\color{muted}As decisões do projeto são justificadas}
  \\item[$\\square$] {\\small\\color{muted}A argumentação é coerente com os documentos de DT e DCU}
\\end{itemize}
\\medskip

\\subsubsection*{Referências e Fontes \\hfill {\\chiplow 0,00 / 0,25 pts}}

\\begin{itemize}[leftmargin=18pt,itemsep=2pt,topsep=4pt,parsep=0pt]
  \\item[$\\square$] {\\small\\color{muted}Dados e afirmações têm fonte citada}
  \\item[$\\square$] {\\small\\color{muted}Há pelo menos 3 fontes referenciadas}
  \\item[$\\square$] {\\small\\color{muted}Referências formatadas de forma consistente}
\\end{itemize}

\\subsection*{Protótipo Final \\hfill {\\small\\color{muted}0,00 / 1,00 pts}}
\\noindent\\rule{\\linewidth}{0.3pt}\\vspace{-4pt}

\\subsubsection*{Protótipo Navegável e Completo \\hfill {\\chiplow 0,00 / 0,50 pts}}

\\begin{itemize}[leftmargin=18pt,itemsep=2pt,topsep=4pt,parsep=0pt]
  \\item[$\\square$] {\\small\\color{muted}O protótipo está navegável no Figma — não são telas soltas}
  \\item[$\\square$] {\\small\\color{muted}Cobre o fluxo principal com pelo menos 3 telas conectadas}
  \\item[$\\square$] {\\small\\color{muted}Um usuário externo consegue navegar sem explicação}
\\end{itemize}
\\medskip

\\subsubsection*{Fidelidade à Identidade Visual \\hfill {\\chiplow 0,00 / 0,50 pts}}

\\begin{itemize}[leftmargin=18pt,itemsep=2pt,topsep=4pt,parsep=0pt]
  \\item[$\\square$] {\\small\\color{muted}O protótipo é consistente com o Style Guide}
  \\item[$\\square$] {\\small\\color{muted}Cor, tipografia e componentes seguem o planejado}
\\end{itemize}

\\subsection*{Pitch + Slides \\hfill {\\small\\color{muted}0,00 / 1,50 pts}}
\\noindent\\rule{\\linewidth}{0.3pt}\\vspace{-4pt}

\\subsubsection*{Clareza e Fluidez da Apresentação \\hfill {\\chiplow 0,00 / 0,50 pts}}

\\begin{itemize}[leftmargin=18pt,itemsep=2pt,topsep=4pt,parsep=0pt]
  \\item[$\\square$] {\\small\\color{muted}O grupo apresenta com segurança — não lê do slide}
  \\item[$\\square$] {\\small\\color{muted}Todos os integrantes participam}
  \\item[$\\square$] {\\small\\color{muted}O tempo de 7 minutos é respeitado}
\\end{itemize}
\\medskip

\\subsubsection*{Slides — Qualidade Visual \\hfill {\\chiplow 0,00 / 0,50 pts}}

\\begin{itemize}[leftmargin=18pt,itemsep=2pt,topsep=4pt,parsep=0pt]
  \\item[$\\square$] {\\small\\color{muted}Os slides são visuais — não são paredes de texto}
  \\item[$\\square$] {\\small\\color{muted}Hierarquia tipográfica clara}
  \\item[$\\square$] {\\small\\color{muted}A identidade visual do projeto está presente nos slides}
\\end{itemize}
\\medskip

\\subsubsection*{Capacidade de Justificar Decisões \\hfill {\\chiplow 0,00 / 0,50 pts}}

\\begin{itemize}[leftmargin=18pt,itemsep=2pt,topsep=4pt,parsep=0pt]
  \\item[$\\square$] {\\small\\color{muted}O grupo responde perguntas sobre as escolhas com segurança}
  \\item[$\\square$] {\\small\\color{muted}Sabe explicar por que fez o que fez}
  \\item[$\\square$] {\\small\\color{muted}Quando não sabe, reconhece honestamente}
\\end{itemize}

\\subsection*{Processo Ágil — Trello \\hfill {\\small\\color{muted}0,00 / 1,50 pts}}
\\noindent\\rule{\\linewidth}{0.3pt}\\vspace{-4pt}

\\subsubsection*{Organização e Uso Real do Trello \\hfill {\\chiplow 0,00 / 0,75 pts}}

\\begin{itemize}[leftmargin=18pt,itemsep=2pt,topsep=4pt,parsep=0pt]
  \\item[$\\square$] {\\small\\color{muted}O quadro tem as colunas corretas: Backlog, To Do, Doing, Done e Referências}
  \\item[$\\square$] {\\small\\color{muted}Os cards têm título no formato verbo no imperativo}
  \\item[$\\square$] {\\small\\color{muted}O histórico mostra uso ao longo do semestre}
\\end{itemize}
\\medskip

\\subsubsection*{Sprint Reviews \\hfill {\\chiplow 0,00 / 0,75 pts}}

\\begin{itemize}[leftmargin=18pt,itemsep=2pt,topsep=4pt,parsep=0pt]
  \\item[$\\square$] {\\small\\color{muted}Ao menos 3 Sprint Reviews foram preenchidos}
  \\item[$\\square$] {\\small\\color{muted}Os reviews têm reflexão real: entregas, pendências, impedimentos e aprendizados}
  \\item[$\\square$] {\\small\\color{muted}Há honestidade sobre o processo}
\\end{itemize}

\\subsection*{Organização e Nomenclatura de Arquivos \\hfill {\\small\\color{muted}0,00 / 0,50 pts}}
\\noindent\\rule{\\linewidth}{0.3pt}\\vspace{-4pt}

\\subsubsection*{Estrutura e Padrão de Nomenclatura \\hfill {\\chiplow 0,00 / 0,50 pts}}

\\begin{itemize}[leftmargin=18pt,itemsep=2pt,topsep=4pt,parsep=0pt]
  \\item[$\\square$] {\\small\\color{muted}Todos os arquivos estão na pasta correta do Drive}
  \\item[$\\square$] {\\small\\color{muted}Os arquivos seguem o padrão de nomenclatura}
  \\item[$\\square$] {\\small\\color{muted}A indicação de turma (\_A ou \_B) está correta em todos}
\\end{itemize}

\end{document}
